\subsection{The canonical divisor and the exact anticanonical ring}
\label{sec:canonical-ring-recovery}

The analytic construction admits a further explicit calculation which does
not use its claimed identification with a sphere. Retain the actual proper
map $f:X\to B=\mathbb P^1$, the cusp $p_0$, the finite points $p_1,p_2$,
the reduced finite fibres $S_1,S_2$, and
$f^*(p_1)=3S_1$, $f^*(p_2)=4S_2$. All line-bundle isomorphisms below
refer to this constructed analytic space. No replacement period matrix or
new compactification is introduced.

\begin{theorem}[Canonical bundle and recovery of the original fibration]
\label{thm:canonical-ring-recovery}
The evaluation map
\[
 f^*(f_*\omega_{X/B})\longrightarrow\omega_{X/B}
\]
has zero divisor exactly $2S_2$. Consequently
\[
 \omega_X\simeq f^*\mathcal O_B(-p_2)\otimes\mathcal O_X(2S_2)
       \simeq\mathcal O_X(-2S_2),\qquad
 \omega_X^{-2}\simeq f^*\mathcal O_B(p_2).
 \tag{CR1}
\]
For every integer $k$ the actual direct image of the vertical divisor is
\[
 f_*\mathcal O_X(kS_2)
       =\mathcal O_B\bigl(\lfloor k/4\rfloor p_2\bigr)
 \tag{CR2}
\]
as subsheaves of the meromorphic-function sheaf, using pullback to identify
their sections. In particular, for each integer $m\geq0$,
\[
 h^0(X,\omega_X^{-m})=\lfloor m/2\rfloor+1,
 \qquad h^0(X,\omega_X^m)=0\quad(m>0).
 \tag{CR3}
\]
There is a graded algebra isomorphism
\[
 \mathbb C[U,V]\xrightarrow{\sim}
       \bigoplus_{m\geq0}H^0(X,\omega_X^{-m}),
       \qquad\deg U=1,\quad\deg V=2.
 \tag{CR4}
\]
The degree-two complete linear system has no base points and its map is
the original $f$, followed by a specified projective-coordinate map of $B$.
The degree-one system consists of the single divisor $2S_2$.
\end{theorem}

\begin{proof}
Write $\mathfrak K=f_*\omega_{X/B}$. The map is the natural adjunction
evaluation, not an identification inferred from the ranks of its two
terms. On a smooth torus fibre the invariant top form is nowhere zero.
At the cusp, the original form
$(dx_1/x_1)\wedge(dx_2/x_2)\wedge dt_c$ is a nowhere-zero canonical
frame in each of the unimodular toric charts. Its relative version is
the local generator of $\mathfrak K$ used in
\cref{lem:hodge-line-audit}, so evaluation has no zero there.

For a finite filling retain $t_j=s_j^{m_j}$,
$(m_1,m_2)=(3,4)$ and $(a_1,a_2)=(2,1)$. On its finite \'etale
total-space cover $q_j:J'_j\to N_j$, the actual direct-image generator
pulls back, by (9.5c)--(9.5f), to
\[
 s_j^{m_j-1-a_j}e^{\varphi_j(s_j)}
       (ds_j\wedge e)\otimes(dt_j)^{-1}.
 \tag{CR5}
\]
Here $(dt_j)^{-1}$ is the pulled-back inverse canonical frame on the
base, not division by $m_js_j^{m_j-1}ds_j$. The coefficient
$e^{\varphi_j}$ is a unit. Thus the zero divisor on the cover is
$(m_j-1-a_j)\{s_j=0\}$. Since $q_j$ is \'etale, faithful pullback
of Cartier divisors makes the downstairs divisor
$(m_j-1-a_j)S_j$. The two coefficients are respectively $0$ and $2$.
These neighbourhoods and the smooth locus cover $X$, proving the stated
global zero divisor with no omitted component.

The evaluation therefore identifies
\[
 \omega_{X/B}\simeq f^*\mathfrak K\otimes\mathcal O_X(2S_2).
\]
By \cref{lem:hodge-line-audit}, $\mathfrak K\simeq\mathcal O_B(p_0)$.
In the original affine base coordinate $t$, with
$t(p_1)=0$, $t(p_2)=1$, $t(p_0)=\infty$, retain the explicit choices
\[
 h=\frac1{t-1},\qquad \eta_B=\frac{dt}{(t-1)^2}.
\]
Their divisors are $p_0-p_2$ and $-2p_2$, respectively: at infinity,
with $w=1/t$, they are $w/(1-w)$ and $-dw/(1-w)^2$.
Thus $\eta_B$ specifies $\omega_B\simeq\mathcal O_B(-2p_2)$.
Tensoring with $f^*\omega_B$ gives
\[
 \omega_X\simeq
 \mathcal O_X\bigl(f^*(p_0)-2f^*(p_2)+2S_2\bigr).
\]
Multiplication by $f^*h$, with its displayed divisor, identifies the
base factor $\mathcal O_B(p_0-2p_2)$ with $\mathcal O_B(-p_2)$.
Using the actual multiplicity $f^*(p_2)=4S_2$ proves (CR1).

To prove (CR2), a section of $\mathcal O_X(kS_2)$ is represented by
a meromorphic function on $X$ with the indicated vertical order bound.
Away from $p_2$ it is holomorphic. Its restriction to each connected
compact fibre there is constant; $f_*\mathcal O_X=\mathcal O_B$
therefore descends it uniquely to a holomorphic function on that base
open set. Near $p_2$, pull back to the original cover $J'_2$, including
its zero section. The descended function becomes a function $a(s_2)$,
and $s_2^k a(s_2)$ is holomorphic at zero (with the evident vanishing
interpretation when $k<0$). Thus $a$ is meromorphic and has a Laurent
series bounded below. Invariance under the actual quotient, whose base
coordinate is $s_2\mapsto\zeta_2s_2$, says
$a(\zeta_2s_2)=a(s_2)$. Only exponents divisible by four occur, so
$a(s_2)=b(s_2^4)$ for a meromorphic $b$ on the original base disc.
The order inequality is exactly
\[
 4\operatorname{ord}_{p_2}(b)+k\geq0
 \quad\Longleftrightarrow\quad
 \operatorname{ord}_{p_2}(b)\geq-\lfloor k/4\rfloor.
\]
Conversely every such $b$ pulls back to an allowed meromorphic function.
This proves equality in both directions, as well as compatibility on
overlaps. No assumption about the field of all meromorphic functions
on $X$ is used. Applying (CR2) with $k=2m$ and $k=-2m$ and using
$h^0(\mathbb P^1,\mathcal O(d))=\max\{d+1,0\}$ proves (CR3).

Here is the exact algebra map in (CR4). Fix the isomorphism
$\omega_X^{-1}\simeq\mathcal O_X(2S_2)$ obtained above. Let $u$ be
the section corresponding to the canonical divisor section of
$\mathcal O_X(2S_2)$, so $\operatorname{div}(u)=2S_2$.
Its square, under the induced degree-two identification in (CR1),
is $f^*\sigma_{p_2}$, where $\sigma_{p_2}$ is the canonical divisor
section of $\mathcal O_B(p_2)$. Choose a second section $\ell$ of this
line bundle with $\ell(p_2)\ne0$, retaining that choice, and let $v$
correspond to $f^*\ell$ under the same degree-two identification.
Define $U\mapsto u$, $V\mapsto v$.

For $m=2r$, the monomials $u^{2r-2j}v^j$, $0\leq j\leq r$,
correspond exactly to
$f^*(\sigma_{p_2}^{r-j}\ell^j)$, a basis of
$H^0(B,\mathcal O_B(rp_2))$. For $m=2r+1$, multiplication by $u$
identifies the same basis with the full space: in the meromorphic
description (CR2), both $4rS_2$ and $(4r+2)S_2$ allow precisely
the base pole orders at most $r$. Thus the displayed map is an
isomorphism in every degree, proving both surjectivity and absence of
polynomial relations. Replacing $\ell$ by $a\ell+b\sigma_{p_2}$,
$a\ne0$, changes the second generator by exactly
$v\mapsto av+bu^2$; no scalar is silently discarded.
Replacing the chosen isomorphism $\omega_X^{-1}\to\mathcal O_X(2S_2)$
by $c$ times itself changes the generators by
$(u,v)\mapsto(c^{-1}u,c^{-2}v)$, with $c\in\mathbb C^*$ retained.

The two degree-two generators have no common zero because
$\sigma_{p_2},\ell$ do not. Their map is
$x\mapsto[\sigma_{p_2}(f(x)):\ell(f(x))]$, which is exactly
$f$ followed by that chosen coordinate map $B\to\mathbb P^1$.
Explicitly, writing $\ell=A/(t-1)+B$ with $A\ne0$, it is
$t\mapsto[t-1:A+B(t-1)]$, including the values $[0:A]$ at $p_2$
and $[1:B]$ at $p_0$.
Degree one has basis $u$ and its divisor is $2S_2$, as claimed.
\end{proof}

In the already specified additive Picard coordinate (9.19), with
$L_{12}=f^*\mathcal O_B(1)$, the identities retain the original constants:
\[
 \mathcal O_X(S_1)=L_4,\qquad
 \mathcal O_X(S_2)=L_3,\qquad
 \omega_X=L_{-6},\qquad \omega_X^{-2}=L_{12}.
\]
This last expression uses the previously proved Picard identification;
the divisor, direct-image and ring proofs above do not depend on it.
